\documentclass{article}

\textwidth=6.5in
\textheight=8.5in
\voffset=-54pt
\oddsidemargin=0in
\evensidemargin=0in
\setlength{\parskip}{8pt}
\setlength{\parindent}{0pt}

\usepackage{amsmath, amssymb}
\usepackage{ifthen}

\usepackage[inline]{enumitem}
\setlist{topsep=0pt, parsep=8pt}

\usepackage[rgb,table]{xcolor}\convertcolorsUfalse
\usepackage{graphicx}

% change hyperref options
\PassOptionsToPackage{hyphens}{url}
\usepackage[bookmarks,colorlinks,linkcolor=black,citecolor=blue,pdfstartview=FitH,urlcolor=blue]{hyperref}
\hypersetup{pdfpagemode=UseNone}
\usepackage{bookmark}

\usepackage{graphicx}
\usepackage{multicol}
\usepackage{framed}
\usepackage{array}

\usepackage{icomma}

\usepackage{soul}
\usepackage[normalem]{ulem}

\usepackage{pgfplots}
\pgfplotsset{compat=1.18}  
\pgfplotscreateplotcyclelist{mycolor}{black, blue, red, green!70!black, magenta, purple, cyan, orange }  
\pgfplotsset{
  disabledatascaling,
  cycle list=tikzcolor, 
  axis lines=middle,
  every axis plot/.style={no marks, thick},
  every axis/.style={
    enlarge x limits=false,
    enlarge y limits=auto,
    xlabel style={at={(current axis.right of origin)},anchor=west},
    ylabel style={at={(current axis.above origin)},anchor=south},
    % extra x and y ticks for asymptotes
    extra x tick style={grid=major, grid style={dashed,black}},
    extra y tick style={grid=major, grid style={dashed,black}},
    /pgf/number format/1000 sep={},
  },    
  tick label style = {font=\footnotesize, fill=\currentbgcolor, inner sep=0pt},    
  xticklabel shift=1pt,    
  scaled ticks=false,
  scaled y ticks=false,
  unbounded coords=jump,
  samples=101,
  minor tick length=0,
  scatter/classes={
    open={mark=*,draw=black,fill=white, mark size=2, thin},
    closed={mark=*,draw=black,fill=black, mark size=2, thin}
  },
  width=3in,
}
\pgfplotsset{open/.style={only marks, mark=*, fill=white, thin, mark size=2}}
\pgfplotsset{closed/.style={only marks, mark=*, draw=black, fill=black,
    thin, mark size=2}}
\pgfplotsset{labeled/.style={nodes near coords, only marks, mark=*, draw=black, fill=black, thin, mark size=2, point meta=explicit symbolic}}

\usepackage{tikz}
\usetikzlibrary{
  matrix,%
  % enables the snake paths
  decorations.pathmorphing,%
  % enables bigger arrows
  decorations.markings,%
  decorations.pathreplacing,%
  decorations.text,%
  calc,%
  patterns,%
  shapes.geometric,%
  arrows,
  decorations.shapes,
  positioning,
  plotmarks,
  intersections
}
\tikzset{>=latex} % sets default arrow type in tikz
\tikzset{font=\normalsize}
\tikzset{mark size=1.5pt, mark options=thin}
\tikzset{pin distance=4pt,
  every pin edge/.style={<-, thin, shorten <= -2pt}}

\interfootnotelinepenalty=10000
\renewcommand{\thefootnote}{(\arabic{footnote})}

\usepackage{multirow, bigdelim}

% change to \usepackage[sols]{handout} to see solutions
\usepackage[sols]{handout}

\newcommand\iscurrentenumi[1]{%
  \ifthenelse{\equal{#1}{\theenumi}}{}{#1}}
\setenumerate[1]{ref=\#\arabic*}
\setenumerate[2]{ref=\protect\iscurrentenumi{\theenumi}(\alph*)}
\newcommand\enumiiprefix[2]{%
  \ifthenelse{{\equal{#1}{\theenumi}}\AND{\equal{#2}{(\alph{enumii})}}}{}{}%
  \ifthenelse{{\equal{#1}{\theenumi}}\AND{\NOT{\equal{#2}{(\alph{enumii})}}}}{#2}{}%
  \ifthenelse{\equal{#1}{\theenumi}}{}{#1#2}%
}
\setenumerate[3]{ref=\protect\enumiiprefix{\theenumi}{(\alph{enumii})}\roman*}

\newlength{\tipmargin}
\makeatletter

\renewenvironment{framed}{
  \setlength{\tipmargin}{\textwidth-\linewidth}
  \advance\@totalleftmargin-\tipmargin
  \def\FrameCommand{\hspace{\tipmargin}\fbox}%
  \advance\linewidth\tipmargin
  \MakeFramed {\advance\hsize-\width \FrameRestore}%
}
{\endMakeFramed}

\makeatother

\definecolor{purple}{rgb}{0.5,0,1.0}

\usepackage{fancyhdr}
\lhead{\textcolor{white}{This content is protected and may not be shared, uploaded, or distributed.}}
\rhead{}
\rfoot{\vspace*{\fill}\footnotesize\textcopyright{} \emph{Harvard Math Ma}}
\renewcommand{\headrulewidth}{0pt}
\pagestyle{fancy}
\begin{document}

\htitle{34. Quadratics}

\begin{problemonly}
  \begin{framed}

You should be able to:
    \begin{itemize}[topsep=0pt,partopsep=0pt,parsep=8pt,itemsep=0pt]
    \item Explain why, if acceleration is constant, then position is a quadratic function of time.

    \item Graph and analyze quadratics of the forms $f(x) = ax^2 + bx + c$, $f(x) = a (x - r_1)(x - r_2)$, and $f(x) = a (x - h)^2 + k$.

    \item Write a quadratic that fits given information, choosing the most appropriate form.

    \end{itemize}
  \end{framed}
\end{problemonly}

\begin{enumerate}
\item  \label{gravity-apple}

  One day, Katie leans out her dorm window and hurls an apple straight up in the air.

  \begin{enumerate}

  \item
    \begin{problem}[\vfill]
      We'd like to model $h(t)$, the height of the apple in feet $t$ seconds after Katie tosses it. Sketch a plausible graph of $h(t)$.
    \end{problem}

    \begin{solution}
      \begin{center}
        \begin{tikzpicture}
          \begin{axis}[xlabel=$t$,ylabel=$h$,xtick=\empty, ytick=\empty]
            \addplot[domain=0:4] {-16*x^2+48*x+64};
          \end{axis}
        \end{tikzpicture}
      \end{center}
    \end{solution}

  \item

    \begin{problem}[\stretch{0.35}]
      Does it seem reasonable to model $h(t)$ with a linear function? An exponential function? A logarithmic function?
    \end{problem}

    \begin{solution}
      None of these seem reasonable. Linear functions,
      exponential functions, and logarithmic functions are either
      always increasing or always decreasing, whereas $h(t)$
      increases for a bit then decreases.
    \end{solution}

  \item

    \note{I want students to identify that, once Katie lets go of the apple, the only thing affecting its motion is gravity. (Well, there's also air resistance, but we'll ignore that.)

      I'll let students know that gravity produces a constant acceleration of $-32$ ft/s$^2$. This isn't very intuitive if you've never encountered it before! But you can help students make sense of it; for instance, if you drop a ball from 6 feet up and another from 60 feet up, you'd expect the second one to be moving a lot faster when it hits the ground; that's because it's been accelerating.}

    \begin{problem}[\stretch{0.375}]
      What do we know about $h(t)$? Express what we know in function notation. 
    \end{problem}    

    \begin{solution}
      Once Katie lets go of the apple, the only thing causing it to move is gravity. Gravity causes the apple to have a constant acceleration of $-32$ ft/s$^2$. That is, if we let $a(t)$ be the apple's acceleration $t$ seconds after Katie lets go, then $a(t) = -32$.
    \end{solution}

    \note{We need to know the ball's initial height and initial velocity to have any hope of knowing where it is at any time $t$.}

    \begin{problem}[\nstretch{0.375}]
      Do we have enough information to find a formula for $h(t)$? If not, what additional information do we need?
    \end{problem}

    \begin{solution}      
      To know what happens to the apple, we also need to know how hard Katie threw it, which can be expressed by giving the initial velocity, and how high up Katie threw the apple from, which is the initial height.      
    \end{solution}

  \end{enumerate}

  Suppose that Katie's window is 64 ft above the ground, and she hurls the apple with an initial velocity of 48 ft/s.

  \begin{enumerate}[resume]
  \item
    \begin{problem}[0.25in]
      Express this information in function notation. 
    \end{problem}

    \begin{solution}
      If Katie's window is 64 feet above the ground, then $h(0)=64$.
      If the apple's initial upward velocity is $48$ ft/s and
      we use $v(t)$ for the apple's velocity at time $t$, then $v(0)=48$.
    \end{solution}

  \item \label{gravity-apple-height}
    \begin{problem}[\vfill]
      Find a formula for $h(t)$, the apple's height in feet $t$ seconds after Katie tosses it.
    \end{problem}

    \begin{solution}
      Since acceleration is the derivative of velocity, we have $v'(t) = -32$. We know many functions with derivative $-32$: a few examples are $-32 t$, $-32t + 1$, and $-32 t - 700$. We can generalize this by saying that the functions with derivative $-32$ are those of the form $v(t) = -32 t + C$ where $C$ is a constant. But we're trying to think of an example that has $v(0) = 48$, so we must have $v(t) = -32 t + 48$.

      We can try the same idea again; velocity is the derivative of position, so $h'(t) = -32 t + 48$. The functions with derivative $-32 t + 48$ are $h(t) = - 16 t^2 + 48 t + C$ where $C$ is a constant. Since we want $h(0)$, we must have $h(t) = -16 t^2 + 48 t + 64$.
    \end{solution} 

  \item \label{gravity-apple-max}
    \begin{problem}[\stretch{0.35}]
      When does the apple reach its greatest height? 
    \end{problem}

    \begin{solution}
      The apple reaches its greatest height when its velocity changes
      from positive to negative. This happens when $v(t)=0$.
      Since $v(t)=-32t+48$, the apple reaches its maximum height
      at $t=1.5$ seconds.
    \end{solution}

  \item

    \note{Here, I will have students factor $h(t)$.}

    \begin{problem}[\nstretch{0.5}]
      When does the apple hit the ground?
    \end{problem}

    \begin{solution}
      The apple hits the ground when $h(t)=0$. To solve for $t$,
      we can factor $h(t)$:
      \begin{align*}
        h(t) &= -16t^2 + 48t + 64 \\
        &= -16(t^2 - 3t - 4) \\
        &= -16(t-4)(t+1)
      \end{align*}
      The solutions to $h(t)=0$ are $t=-1$ and $t=4$. This means that the 
      apple hits the ground after $4$ seconds. (The solution $t=-1$ is
      extraneous because the domain of $h(t)$ is $t\geq 0$.)
    \end{solution}

  \item

    \note{Here, I want them to sketch the entire quadratic (even though $h(t)$ only makes sense for $t \geq 0$). I will highlight that, in the factored form, it's easy to read off the roots. Also, the vertex is half-way between the roots. That means we could have done \ref{gravity-apple-max} in another way: rather than setting $h'(t) = 0$, we also could have just found the roots of $h(t)$ and concluded the maximum is half-way between them.}

    \begin{problem}[\vfill]
      Sketch a graph of the formula you found in \ref{gravity-apple-height}. 
    \end{problem}

    \begin{solution}
      To sketch $y=-16x^2+48x+64$, 
      we know that the $x$-intercepts will be at $x=-1$ and $x=4$, since
      the factored form is $y=-16(x+1)(x-4)$.
      We also know that the $y$-intercept is at $y=64$. The maximum is
      at $x=1.5$, halfway between the $x$-intercepts.
      \begin{center}
        \begin{tikzpicture}
          \begin{axis}[xtick={-1,1.5,4}, ytick={64}]
            \addplot[domain=-2:5] {-16*x^2+48*x+64};
          \end{axis}
        \end{tikzpicture}
      \end{center}
      
    \end{solution}

  \end{enumerate}

\end{enumerate}
\note{In the previous problem, we will have written $h(t)$ in two ways, $h(t) = - 16t^2 + 48 t + 64 = -16 (t + 1)(t - 4)$. We'll generalize this to the two forms $f(x) = ax^2 + bx + c$ and $f(x) = a (x - r_1)(x - r_2)$.}

\begin{framed}
  A function $f$ is \uline{quadratic} if it can be written as $f(x) = ax^2 + bx + c$ for some constants $a, b, c$ with $a \neq 0$. 

  \textbf{3 ways of writing a quadratic.}

  \vspace*{-12pt}
  \renewcommand{\arraystretch}{2}
  \begin{tabular}{l p{1.5in} p{4in}}
    & \emph{Form} & \emph{Easily accessible information...} \\[-4pt]
    \cline{2-3} 
    1. & $f(x) = ax^2 + bx + c$ & \solonly{$y$-intercept}\\
    2. & \solonly{$f(x) = a (x - r_1)(x - r_2)$} & \solonly{roots} \\
    3. & \solonly{$f(x) = a (x - h)^2 + k$} & \solonly{graph, vertex} 
  \end{tabular}
  \renewcommand{\arraystretch}{1}  
\end{framed}

\begin{enumerate}[resume]
\item 
  \begin{problem}
    Which of the following are quadratics?
  \end{problem}

  \begin{multicols}{2}
    \begin{enumerate}[label=\maybecircle{\Alph*.}]
    \circleitem $f(x) = 2x^2 - 7x + 3$
    \plainitem $f(x) = \dfrac{1}{x^2 + 1}$
    \circleitem $f(x) = (3 - x)(5 - 7x) + 4$
    \plainitem $f(x) = e^{x^2}$
    \end{enumerate}
  \end{multicols}

  \pspace{\nstretch{0.25}}

\item  

  \note{After this problem, we'll generalize this to say the third form of a quadratic is $a(x - h)^2 + k$.}

  \begin{enumerate}
  \item \label{sketch-quadratic-1}

    \note{Practice sketching using shifts and stretches.}

    \begin{problem}[\vfill]
      Sketch the graph of $y = -2 (x - 1)^2 + 18$.  (If you start with $y = x^2$, what shifts and stretches do you need to do to get this graph?) Where's its vertex?
    \end{problem}

    \begin{solution} 
      Here the graph is a combination of graph transformations. Remember to follow the order of operations. Translate one unit to the right, stretch by a factor of 2, reflect over the $x$-axis, and then translate up by 18 units.

      \begin{tikzpicture} 
        \begin{axis}[xtick={1}, ytick=18]
          \addplot+[domain=-2:4] {-2*(x-1)*(x-1)+18};
        \end{axis}
      \end{tikzpicture}

      Note that you can quickly read off the vertex of quadratic functions written in this form.
    \end{solution}

  \item

    \begin{problem}[\vfill]
      Does the graph have $x$-intercepts?  If so, where?
    \end{problem}

    \begin{solution} 
      \begin{align*} 
        -2 (x - 1)^2 + 18 &=0 \\  
        -2(x-1)^2 &= -18 \\
        (x-1)^2 & = 9 \\ 
        x-1 & = \pm3 \\ 
        x & = 1\pm3 \\
        &= -2, 4 
      \end{align*} 
    \end{solution} 

  \item
    \begin{problem}[\vfill\newpage]
      How would you rewrite $y = - 2 (x - 1)^2 + 18$ in the other 2 ways? 
    \end{problem}

    \begin{solution}
      To rewrite $y=-2(x-1)^2+18$ in the first form, we can simply
      expand $(x-1)^2$ to get $y=-2(x^2-2x+1)+18$, then distribute
      and combine like terms, $y=-2x^2+4x+16$.

      To rewrite into factored form, we can take $y=-2x^2+4x+16$ and
      factor it.
      \begin{align*}
        y &= -2x^2 + 4x + 16 \\
        &= -2(x^2-2x-8) \\
        &= -2(x+2)(x-4)
      \end{align*}
      We could've also found this without factoring by noticing
      that $a=-2$ and using the information that the $x$-intercepts
      are $-2$ and $4$.
    \end{solution}

  \end{enumerate}

\item 

  \note{Here, I think it will come up that $a$ tells us which way the parabola opens. If it doesn't, I will bring it up!}

  In each part, you're given a quadratic. Sketch its graph, and find the coordinates of its vertex. Does the function have a global maximum on $(-\infty, \infty)$? How about a global minimum? 

  \begin{enumerate}
  \item
    \begin{problem}[\vfill]
      $f(x) = 2 (x - 1) (x - 5)$
    \end{problem}

    \begin{solution}
      We can tell that, if we multiplied this out, the coefficient of $x^2$ would be $2$, so the parabola opens upward. Also, it has $x$-intercepts at $x = 1$ and $x = 5$. By symmetry, the vertex of the parabola is half-way between the two $x$-intercepts, so it's at $x = 3$. If we plug $x = 3$ into function, we get $f(3) = 2 \cdot 2 \cdot (-2) = -8$, so the vertex is at $(3, -8)$. This is the global minimum, and there is no global maximum.
      \begin{center}
        \begin{tikzpicture}
          \begin{axis}[xtick={1, 3, 5}, ytick={-8}, width=2in]
            \addplot+[domain=-1:7] {2*(x-1)*(x-5)};
          \end{axis}
        \end{tikzpicture}
      \end{center}
    \end{solution}

  \item
    \begin{problem}[\vfill]
      $g(x) = -2x^2 + 4x - 7$
    \end{problem}

    \begin{solution}
      The vertex happens when $g'(x)=0$. Since $g'(x) = -4x+4$, this happens at $x = 1$. So, the vertex is at $(1, g(1)) = (1, -5)$.
      Since the coefficient of $x^2$ is negative, the parabola opens
      downwards, so this is the global maximum. There is no global minimum.
      \begin{center}
        \begin{tikzpicture}
          \begin{axis}[xtick={1}, ytick={-5}, ymax=5, width=2in]
            \addplot+[domain=-1.5:3.5] {-2*x*x+4*x-7};
          \end{axis}
        \end{tikzpicture}
      \end{center}
    \end{solution}

    \note{The next question is to make the point that this quadratic has no real roots, so we can't write $g(x)$ in the form $a(x - r_1)(x - r_2)$ (at least not with $r_1, r_2$ being real numbers).}

    \begin{problem}[\stretch{0.35}]
      Can you rewrite $g(x)$ in the other 2 ways?
    \end{problem}

    \begin{solution}
      We can rewrite $g(x)$ in vertex form with 
      $a=-2$, $h=1$, and $k=-5$. So $g(x)=-2(x-1)^2-5$. However,
      we can see that the graph of $y=g(x)$ has no $x$-intercepts,
      which means that we cannot rewrite $g(x)$ in factored form
      using \textit{real numbers} $r_1$ and $r_2$. 
    \end{solution}
  \end{enumerate}

  \probonly{\emph{Even though this question was the same as \ref{sketch-quadratic-1}, how did your strategy differ?}}

\item  In each part, how many quadratics are there with the given properties? Can you write the equation of one such quadratic?

  \begin{enumerate}
  \item

    \note{I hope we'll have time to do this just to give students an example of the strategy: we can write an equation $y = a(x - 1)(x - 4)$ with a parameter and then solve for that parameter.}

    \begin{problem}[\nstretch{0.65}]
      A parabola with $x$-intercepts at $x = 1$ and $x = 4$ and $y$-intercept $-1$. 
    \end{problem}

    \begin{solution}
      Based on the $x$-intercepts, we know the parabola will have the form $y = a (x - 1) (x - 4)$. We also know $(0, -1)$ is a point on the parabola, so $-1 = a (0-1)(0-4)$. Therefore, $a = - \frac{1}{4}$. Thus, $\boxed{y =  - \frac{1}{4} (x - 1) (x - 4)}$, and there is only \fbox{one such parabola}.
    \end{solution}

  \item
    \begin{problem}[\vfill]
      A quadratic function $f$ such that $f(3) = 5$, $f'(3) = 0$, and $f''(3) > 0$. 
    \end{problem}

    \begin{solution}
      Since $f'(3)=0$, the vertex must occur at $x=3$. We know
      that $f(3)=5$, so the vertex is $(3, 5)$. This means
      that $f(x)=a(x-3)^2+5$. Since $f''(3)>0$, the parabola
      must be concave up, so $a>0$, but we don't have any
      more information about $a$. So one such function is
      $\boxed{f(x)=(x-3)^2+5}$, and there are \fbox{infinitely
      many} functions that satisfy these properties.
    \end{solution}

  \item

    \note{Since this has only one $x$-intercept, that must be the vertex.}

    \begin{problem}[\vfill]
      A parabola whose $x$- and $y$-intercepts are $(-3, 0)$ and $(0, 6)$. 
    \end{problem}

    \begin{solution}
      Since there is only one $x$-intercept, $(-3, 0)$, it must
      also be the vertex. So $y=a(x+3)^2$. Since the graph
      passes through the point $(0,6)$, $6=a(0+3)^2$, so $a=\frac{2}{3}$.
      This means that the equation of the parabola is 
      $\boxed{y=\frac{2}{3}(x+3)^2}$, and
      there is only \fbox{one such parabola}.
    \end{solution}

  \item

    \note{The easiest strategy here is to start by finding a formula for $f'(x)$.}

    \begin{problem}[\vfill]
      A quadratic function $f$ with $f'(2) = 0$ and $f'(3) = 4$.
    \end{problem}

    \begin{solution}
      The derivative of $f$ is a linear function, and since
      we're given two points for $f'$, we can determine that
      $f'(x)=4x-8$. Based on this, $f(x)=2x^2-8x+c$, where $c$
      is a constant. So one such function is $\boxed{f(x)=2x^2-8x}$,
      and there are \fbox{infinitely many} such functions.
    \end{solution}

  \item

    \note{Students will need to solve a system of equations here. We haven't done this kind of algebra, so it may not be too familiar. This isn't algebra we're going to do again, so don't worry about fluency with this. If students get here, I will encourage them to pick one equation, solve for one variable in it, and plug that into the other equation.}

    \begin{problem}[\vfill]
      A parabola passing through the points $(0, 3)$, $(-1, 6)$, and $(2,
      9)$.
    \end{problem}

    \begin{solution}
      If we write our quadratic in the form $y = ax^2 + bx + c$, then we know
      \begin{align*}
        3 &= c \\
        6 &= a - b + c \\
        9 &= 4a + 2b + c
      \end{align*}
      Substituting $c=3$ into the other two equations, we get the
      system of two equations
      \begin{align*}
        3 &= a-b \\
        6 &= 4a+2b
      \end{align*}
      We can solve the first equation for $a$ to get $a=b+3$.
      Substituting this into the second equation, we get
      \begin{align*}
          3 &= 4(b+3)+2b \\
          6 &= 6b + 12 \\
          -1 &= b
      \end{align*}
      Substituting $b=-1$ into $a=b+3$, we find that $a=2$.
      So the equation of the parabola is $\boxed{y=2x^2-x+3}$, and
      there is only \fbox{one such parabola}.
    \end{solution}

  \end{enumerate}

\end{enumerate}
\end{document}
